\documentclass[11pt]{article}
\usepackage[margin=1in]{geometry}
\usepackage[T1]{fontenc}
\usepackage[utf8]{inputenc}
\usepackage{amsmath,amssymb,amsthm}
\usepackage{enumitem}

\title{ICPC World Finals 2025\\F. Herding Cats}
\author{}
\date{}

\begin{document}
\maketitle

\section*{Problem Summary}

We must permute the $m$ plants so that each cat stops at its prescribed pot. A cat stops at the first
position whose plant belongs to its liked set.

\section*{Key Observations}

\begin{itemize}[leftmargin=*]
    \item For a plant variety $v$, define
    \[
    LB[v] = \max \{p_i : \text{cat } i \text{ likes } v\}.
    \]
    If we placed $v$ before $LB[v]$, then some cat that likes $v$ would stop too early. So $v$ may only
    appear at position $LB[v]$ or later.
    \item Therefore a necessary prefix condition is:
    among the first $j$ pots, we must have at least $j$ plant varieties with $LB[v] \le j$.
    \item Now look at a fixed position $j$ with one or more cats assigned to stop there.
    The plant placed at pot $j$ must be liked by \emph{all} of those cats, and it must satisfy $LB[v]=j$.
    If $LB[v] < j$, then some cat that likes $v$ would have been able to stop earlier than $j$.
    If $LB[v] > j$, then that plant cannot legally be placed at $j$ at all.
    \item These two conditions are also sufficient:
    choose one common plant with $LB=j$ for every occupied position $j$, then place all remaining plants
    in the still-empty positions in nondecreasing order of $LB$.
\end{itemize}

\section*{Algorithm}

\begin{enumerate}[leftmargin=*]
    \item Read all cats and compute $LB[v]$ for every plant $v$.
    \item Check the prefix condition by counting how many plants have each lower bound and accumulating
    these counts from left to right.
    \item Group cats by their target position.
    For each occupied position $j$, count how many cats at $j$ like each plant.
    We need at least one plant that is liked by all cats of that group and also satisfies $LB[v]=j$.
    \item If both checks pass, output \texttt{yes}; otherwise output \texttt{no}.
\end{enumerate}

\section*{Correctness Proof}

We prove that the algorithm returns the correct answer.

\paragraph{Lemma 1.}
In any valid arrangement, every plant $v$ must be placed at a position at least $LB[v]$.

\paragraph{Proof.}
By definition of $LB[v]$, there exists a cat that likes $v$ and is supposed to stop at position $LB[v]$.
If $v$ were placed earlier, that cat would encounter a liked plant before its designated stop and would
stop too soon. \qed

\paragraph{Lemma 2.}
For every occupied position $j$, the plant placed at $j$ must be liked by all cats assigned to $j$ and must
satisfy $LB[v]=j$.

\paragraph{Proof.}
If the plant at $j$ were not liked by some cat assigned to $j$, that cat would continue past pot $j$.
If $LB[v] < j$, then some cat that likes $v$ is supposed to stop earlier, so placing $v$ at $j$ would not be
the first liked plant for all cats that need to stop at $j$.
If $LB[v] > j$, Lemma 1 forbids putting $v$ at $j$. Thus the condition is necessary. \qed

\paragraph{Lemma 3.}
If the prefix condition and the position-wise common-plant condition both hold, then a valid arrangement
exists.

\paragraph{Proof.}
Fix one valid common plant with $LB=j$ for every occupied position $j$.
These choices already guarantee that every such position stops exactly the cats assigned to it.
Now consider the remaining plants and the still-empty positions.
By the prefix condition, for every prefix of length $j$ there are at least $j$ plants with $LB \le j$.
Since we already used exactly one plant with $LB=j$ for each occupied position $j$, the same prefix
condition still holds for the leftovers.
Therefore placing the remaining plants in nondecreasing order of $LB$ fills all empty positions without
ever violating a lower bound. By Lemma 1 no cat stops too early, and the chosen special plants make the
cats at occupied positions stop exactly where required. \qed

\paragraph{Theorem.}
The algorithm outputs the correct answer for every valid input.

\paragraph{Proof.}
If the program answers \texttt{yes}, then both checks pass, and Lemma 3 gives a valid arrangement.
If the program answers \texttt{no}, then either the prefix condition fails, contradicting Lemma 1, or some
occupied position has no valid common plant, contradicting Lemma 2. Hence no valid arrangement
exists. \qed

\section*{Complexity Analysis}

Let $K$ be the total number of liked-plant entries in the test case. The algorithm runs in $O(m+n+K)$
time and uses $O(m+n+K)$ memory. Over all test cases this matches the input bounds.

\section*{Implementation Notes}

\begin{itemize}[leftmargin=*]
    \item The frequency array for the intersection test is reset only on the plants touched by the current
    group of cats, which keeps the total work linear in the input size.
    \item A plant with $LB[v]=0$ is liked by no cat, so it can be used in any still-empty position.
\end{itemize}

\end{document}
